% TOP_PROBLEM: {"id":30007604,"problem_number":"LOCAL-30007604","title":"Digital deposit runs","category_id":21,"category":{"id":21,"name":"economics","display_name":"Economics","description":"Open economic theory statements, published research agendas, and proposed scoped specifications; question provenance and historical priority are qualified per record.","slug":"economics","order_index":21,"created_at":"2026-10-08T00:00:00Z"},"status":"research_question","external_url":null,"legacy_ids":[],"published":false,"statement":"In the explicitly specified setting: How do digital withdrawals and social networks change run thresholds, contagion and the effectiveness of official communication? The mathematical statement above fixes the target, assumptions and scope of this catalogue formulation.","economics":{"original_id":93,"field":"Banking, credit and financial stability","kind":"identification","formalization_basis":"adapted_research_question","source_status":"research_agenda","priority_status":"earliest_found","priority_note":"This is the earliest explicit statement verified in this audit; first priority for the full catalogue formulation is not established.","earliest_verified_reference":{"authors":["Bank of England"],"title":"Bank of England Agenda for Research, 2025–2028","year":2026,"url":"https://www.bankofengland.co.uk/research/bank-of-england-agenda-for-research","doi":null,"locator":"Prudential Architecture / Resilience, paragraph 3","evidence_summary":"Requests technology effects on runs and contagion."},"current_reference":{"authors":["Bank of England"],"title":"Bank of England Agenda for Research, 2025–2028","year":2026,"url":"https://www.bankofengland.co.uk/research/bank-of-england-agenda-for-research","doi":null,"locator":"Prudential Architecture / Resilience, paragraph 3","evidence_summary":"Requests technology effects on runs and contagion."},"formalization_note":"Adapts an explicitly verified research-agenda passage. Mathematical objects, interventions, objectives and scope are proposed here; existing model-specific solutions are not relabeled unsolved."},"statusQualification":"Published question or research agenda verified at the cited locator. The mathematical specification is adapted for this catalogue and includes additional assumptions; source publication does not establish first-ever priority or certify this exact model remains unsolved. Adapts an explicitly verified research-agenda passage. Mathematical objects, interventions, objectives and scope are proposed here; existing model-specific solutions are not relabeled unsolved.","statusReviewedAt":"2026-10-08"}
% FIRST_POSTED: 2026-10-08
% ENTRY_KIND: statement_only
% !TeX program = lualatex
\documentclass[11pt]{article}
\usepackage{fontspec}
\usepackage[a4paper,margin=25mm]{geometry}
\usepackage{amsmath,amssymb,mathtools}
\usepackage{xurl}
\usepackage[hidelinks]{hyperref}
\newcommand{\sref}[2]{\hyperlink{source:#1:#2}{[S#2]}}
\setlength{\emergencystretch}{3em}
\begin{document}
% problemId:30007604
\section*{Digital deposit runs}\label{30007604}
\subsection{Definitions and mathematical statement}
Fix banks \(b=1,\ldots,B\) and a depositor network with observed links or an explicitly modeled latent-network distribution. Each depositor observes signals about bank fundamentals and messages from contacts, then decides when to withdraw subject to a transaction delay \(d\). Let \(v\) denote the speed or reach of social communication and \(m\) an official communication policy. Define \(T_b(d,v,m)\) as the first time cumulative withdrawals exceed a specified fraction \(\rho\in(0,1)\) of initial deposits. The task is to identify the causal effects of lower transaction delays and wider communication on
\[\Pr(T_b\leq h\mid x),\]
for each horizon \(h\) and prespecified fundamental state \(x\), together with contagion to other banks. Separate information transmission from strategic coordination by specifying how signals and messages enter depositor beliefs. Account for endogenous bank pricing, liquid-asset holdings and insurance coverage. Establish exclusions for changes in digital access or communication exposure; uncontrolled simultaneous bad news cannot identify a technology effect. Evaluate which message policies reduce socially costly runs without concealing deteriorating fundamentals. Report heterogeneity across depositor concentration and insurance status, and bounds where the network or signal quality is unobserved.

\par\textbf{Formulation and scope.} Adapts an explicitly verified research-agenda passage. Mathematical objects, interventions, objectives and scope are proposed here; existing model-specific solutions are not relabeled unsolved.

\par\textbf{Open-question provenance.} An explicit question or research agenda was verified at the source locator. This does not by itself certify that every later version remains unresolved \sref{30007604}{1}.

\par\textbf{Historical priority.} This is the earliest explicit statement verified in this audit; first priority for the full catalogue formulation is not established.

\subsection{Short English statement}
In the explicitly specified setting: How do digital withdrawals and social networks change run thresholds, contagion and the effectiveness of official communication? The mathematical statement above fixes the target, assumptions and scope of this catalogue formulation.

\subsection{Sources}
\begin{itemize}\raggedright
\item[\textnormal{[S1]}]\hypertarget{source:30007604:1}{} Bank of England. 2026. Bank of England Agenda for Research, 2025–2028. Prudential Architecture / Resilience, paragraph 3.
\url{https://www.bankofengland.co.uk/research/bank-of-england-agenda-for-research}.
{\small\textit{Earliest verified question/agenda reference; historical priority qualified below. Requests technology effects on runs and contagion.}}
\end{itemize}
\end{document}
